% --- Archon physics formalization source begin ---
% archon:physics
% archon:covers IPhO2026Problems/problem_IPhO_2026_3_A_1.lean
% archon:source-report reports/ipho_2026_k3/problem_IPhO_2026_3_A_1.source.json
% archon:problem-id IPhO_2026_3
% archon:part-id A.1

\chapter{Physics problem IPhO\_2026\_3\_A\_1}
\label{ch:IPhO2026Problems_problem_IPhO_2026_3_A_1}

\paragraph{Problem source.}
A homogeneous isotropic paramagnetic torus has mean radius R, inner radius r
with r << R, volume V, and cross-sectional area A.  An insulated conducting
wire is wound densely around it with N turns and instantaneous current I.
Fields H and B and magnetization M are approximately uniform in the torus.
Use B = mu\_0*H + mu\_0*M, Ampere's law, and the sign convention that work and
heat entering the paramagnetic torus are positive.

Current subquestion:
Write the field magnitude H inside the torus in terms of N, I, A, and V.

\paragraph{Current subquestion.}
Write the field magnitude H inside the torus in terms of N, I, A, and V.

\paragraph{Recorded answer/context.}
H = N*I*A/V.

\paragraph{Figure/image path.}
/root/proposal\_for\_physic/science-mango/ipho\_2026\_source/image/T3\_page-2.png

\paragraph{Formalization target.}
create a compiling Lean file with sorry bodies at `IPhO2026Problems/problem\_IPhO\_2026\_3\_A\_1.lean`.
The Lean declarations must preserve the physical quantities, dimensions or dimensional roles, figure labels, governing-law hypotheses, and final relation expressed by this problem.
Use Mathlib/Physlib names found through LeanExplore where available. If a domain API is missing, introduce faithful local abstractions rather than scalar placeholder aliases.

\begin{theorem}[Physics formalization target]
\label{thm:physics:IPhO_2026_3_A_1:target}
\uses{thm:IPhO2026Problems_problem_IPhO_2026_3_A_1:paramagneticTorus_H_eq_meanRadius}
The assigned autoformalize agent should translate this physics problem into Lean declarations in the covered file, with theorem and lemma proof bodies written as `by sorry`.
\end{theorem}
\begin{proof}
This is an autoformalization task, not a proof task. Produce faithful statements that can later be proved without weakening the source contract.
\end{proof}
% NOTE: PhysLean-coverage exemption (planner-recorded, iter-003): PhysLean's electromagnetism modules have no Ampere-circulation / toroid H-field assembly API for this part's current-operating-point model (see this file's physics-grounding log); the typed `InstantaneousCurrent` amount+dimension wrapper is local by design over the `import Mathlib` baseline. No irrelevant Physlib import is added.

% --- Archon autoformalization helper ledger (added iter-007, coverage-debt batch 2) ---
% Every named Lean declaration of the covered file gets a blueprint entry;
% \uses{} reflects the actual Lean dependency structure read first-hand.

\section*{Declaration ledger (autoformalization contracts)}

\subsection*{Basic habitats and carrier structures}

\begin{definition}[Free space and instantaneous current]
\label{def:IPhO2026Problems_problem_IPhO_2026_3_A_1:FreeSpace}
\lean{IPhO2026.Problem3.PartA1.FreeSpace}
\lean{IPhO2026.Problem3.PartA1.InstantaneousCurrent}
\lean{IPhO2026.Problem3.PartA1.InstantaneousCurrent.ext}
The permeability of free space $\mu_0$ packaged as a named positive record,
and the instantaneous current in the wire as an amount-with-readout wrapper
(its scalar readout in amperes); the extensionality theorem identifies two
currents with equal readouts.
\end{definition}

\begin{definition}[Radial profile and H-field readouts]
\label{def:IPhO2026Problems_problem_IPhO_2026_3_A_1:RadialProfile}
\lean{IPhO2026.Problem3.PartA1.RadialProfile}
\lean{IPhO2026.Problem3.PartA1.HFieldReadouts}
Function-space habitats: a radial profile is a real-valued function of the
radial coordinate, and an H-field readout on a turn-index type is a
per-turn real magnitude readout (A/m).
\end{definition}

\begin{definition}[Amperian filament]
\label{def:IPhO2026Problems_problem_IPhO_2026_3_A_1:AmperianFilament}
\uses{def:IPhO2026Problems_problem_IPhO_2026_3_A_1:FreeSpace}
\lean{IPhO2026.Problem3.PartA1.AmperianFilament}
\lean{IPhO2026.Problem3.PartA1.AmperianFilament.is_free}
\lean{IPhO2026.Problem3.PartA1.AmperianFilament.current_nonneg}
A single turn of the winding treated as an amperian filament: an indexed
turn carrying an instantaneous current certified free and nonnegative;
the consequence theorems re-export freeness and nonnegativity of the
readout.
\end{definition}

\begin{definition}[Ampère circulation law]
\label{def:IPhO2026Problems_problem_IPhO_2026_3_A_1:AmpereLaw}
\lean{IPhO2026.Problem3.PartA1.AmpereLaw}
\lean{IPhO2026.Problem3.PartA1.AmpereLaw.circulation}
\lean{IPhO2026.Problem3.PartA1.AmpereLaw.circulation_constant}
Ampère's law on a finite index of loops: the circulation of the field
magnitude along each loop equals the enclosed free current; the
\texttt{circulation} theorem extracts the per-loop equation and
\texttt{circulation\_constant} its constant-magnitude specialization.
\end{definition}

\begin{definition}[Finite winding]
\label{def:IPhO2026Problems_problem_IPhO_2026_3_A_1:FiniteWinding}
\lean{IPhO2026.Problem3.PartA1.FiniteWinding}
\lean{IPhO2026.Problem3.PartA1.FiniteWinding.card}
A dense winding of $N$ turns: a finite turn-index type with its equivalence
to $\mathrm{Fin}\,N$; \texttt{card} certifies the turn count is exactly $N$.
\end{definition}

\subsection*{Thin-mean-path law and uniformity}

\begin{definition}[Ampère law on the thin mean path]
\label{def:IPhO2026Problems_problem_IPhO_2026_3_A_1:AmpereLawThinMeanPath}
\uses{def:IPhO2026Problems_problem_IPhO_2026_3_A_1:AmpereLaw, def:IPhO2026Problems_problem_IPhO_2026_3_A_1:FiniteWinding}
\lean{IPhO2026.Problem3.PartA1.AmpereLawThinMeanPath}
\lean{IPhO2026.Problem3.PartA1.AmpereLawThinMeanPath.ampere_sum_const}
\lean{IPhO2026.Problem3.PartA1.AmpereLawThinMeanPath.mean_circumference_eq}
\lean{IPhO2026.Problem3.PartA1.AmpereLawThinMeanPath.mean_circumference_mul_eq}
\lean{IPhO2026.Problem3.PartA1.AmpereLawThinMeanPath.mean_circumference_pos}
The thin-torus ($r \ll R$) specialization of Ampère's law: the mean path of
length $2\pi R$ encloses all $N$ turns, each carrying its
current, with per-turn magnitude readout; carries the geometry parameters
$R, r, A, V$ (positive, $V = 2\pi R A$). Consequences: the sum form
$\sum_t 2\pi R\, H_t = \sum_t I_t$, the mean-circumference identities, and
positivity of $2\pi R$.
\end{definition}

\begin{definition}[Uniform field magnitude interface]
\label{def:IPhO2026Problems_problem_IPhO_2026_3_A_1:UniformFieldMag}
\lean{IPhO2026.Problem3.PartA1.UniformFieldMag}
\lean{IPhO2026.Problem3.PartA1.UniformFieldMag.piecewise_eq_H}
\lean{IPhO2026.Problem3.PartA1.UniformFieldMag.uniform_across}
The uniformity interface for an approximately uniform field inside the
torus: a constant magnitude $H$ agreeing with every per-turn readout, so
the piecewise readout equals $H$ and any two turns read the same value.
\end{definition}

\begin{definition}[Uniform-material winding law]
\label{def:IPhO2026Problems_problem_IPhO_2026_3_A_1:AmperianFilamentLaw}
\uses{def:IPhO2026Problems_problem_IPhO_2026_3_A_1:AmpereLawThinMeanPath, def:IPhO2026Problems_problem_IPhO_2026_3_A_1:AmperianFilament, def:IPhO2026Problems_problem_IPhO_2026_3_A_1:UniformFieldMag}
\lean{IPhO2026.Problem3.PartA1.AmperianFilamentLaw}
\lean{IPhO2026.Problem3.PartA1.AmperianFilamentLaw.circulation_eq_filament_current}
\lean{IPhO2026.Problem3.PartA1.AmperianFilamentLaw.filament_is_free}
Along the mean path each turn behaves as an amperian filament carrying the
winding current, and the circulation of the per-turn magnitude equals the
total filament (free) current — the official hint's reading of $I_C$ in
$\oint \vec{H}\cdot d\vec{\ell} = I_C$; the two named theorems are its
elimination and freeness re-export.
\end{definition}

\begin{definition}[Vacuum-core constitutive identity]
\label{def:IPhO2026Problems_problem_IPhO_2026_3_A_1:VacuumCoreIdentity}
\uses{def:IPhO2026Problems_problem_IPhO_2026_3_A_1:FreeSpace}
\lean{IPhO2026.Problem3.PartA1.VacuumCoreIdentity}
\lean{IPhO2026.Problem3.PartA1.VacuumCoreIdentity.b_eq_scaled}
\lean{IPhO2026.Problem3.PartA1.VacuumCoreIdentity.b_uniform}
The electromagnetic material law recorded for downstream parts (not needed
for A.1 itself): pointwise $B = \mu_0 H + \mu_0 M$ with $H$ and $M$
constant-magnitude in the region and $M$ parallel to $H$ (paramagnetic
branch, $M = \chi H$, $\chi \ge 0$); consequences rewrite $B$ as the scaled
field and certify uniformity.
\end{definition}

\subsection*{The paramagnetic torus and the A.1 derivation route}

\begin{definition}[Paramagnetic torus of Part A.1]
\label{def:IPhO2026Problems_problem_IPhO_2026_3_A_1:ParamagneticTorusA1}
\uses{def:IPhO2026Problems_problem_IPhO_2026_3_A_1:FiniteWinding, def:IPhO2026Problems_problem_IPhO_2026_3_A_1:FreeSpace, def:IPhO2026Problems_problem_IPhO_2026_3_A_1:AmpereLawThinMeanPath}
\lean{IPhO2026.Problem3.PartA1.ParamagneticTorusA1}
\lean{IPhO2026.Problem3.PartA1.ParamagneticTorusA1.winding_card}
The full problem-parameter package: $N \ge 1$ turns of dense winding, mean
radius $R > 0$, tube radius $r$ with $0 < r < R$ (thin-torus regime),
cross-section $A > 0$, volume $V = 2\pi R A > 0$, series wire current $I$
(same through every turn), the interior field magnitude $H \ge 0$ under
measurement (never fixed by fiat), and the embedding equations tying the
thin-mean-path Ampère law to these parameters with uniform per-turn
readout $H$. No field states the target $H = NIA/V$.
\end{definition}

\begin{lemma}[Bridge 1 — Ampère in the uniform regime]
\label{lem:IPhO2026Problems_problem_IPhO_2026_3_A_1:ampere_uniform_eq}
\lean{IPhO2026.Problem3.PartA1.ParamagneticTorusA1.ampere_uniform_eq}
\uses{def:IPhO2026Problems_problem_IPhO_2026_3_A_1:ParamagneticTorusA1}
The circulation of the uniform magnitude along the mean path equals the
total threading free current: $2\pi R \cdot (N H) = N I$.
\end{lemma}
\begin{proof}
Sum the per-turn magnitudes (all equal to $H$ by uniformity) against the
per-turn currents (all equal to $I$ by the series winding) in the embedded
thin-mean-path Ampère equation; left to the prover stage.
\end{proof}

\begin{lemma}[Bridge 2 — solve for $H$ along the mean path]
\label{lem:IPhO2026Problems_problem_IPhO_2026_3_A_1:fieldMagnitude_eq_meanRadius_form}
\lean{IPhO2026.Problem3.PartA1.ParamagneticTorusA1.fieldMagnitude_eq_meanRadius_form}
\uses{lem:IPhO2026Problems_problem_IPhO_2026_3_A_1:ampere_uniform_eq}
$H = N I / (2\pi R)$.
\end{lemma}
\begin{proof}
Divide Bridge 1 by $N \cdot 2\pi R \neq 0$ (uses $N \ge 1$, $R > 0$,
$\pi > 0$).
\end{proof}

\begin{lemma}[Bridge 3 — geometry rewrite of the mean circumference]
\label{lem:IPhO2026Problems_problem_IPhO_2026_3_A_1:mean_circumference_eq}
\lean{IPhO2026.Problem3.PartA1.ParamagneticTorusA1.mean_circumference_eq}
\uses{def:IPhO2026Problems_problem_IPhO_2026_3_A_1:ParamagneticTorusA1}
$2\pi R = V / A$.
\end{lemma}
\begin{proof}
Field-simplify the thin-torus volume law $V = 2\pi R A$ using $A > 0$.
\end{proof}

\begin{lemma}[Bridge 4 — figure parametrization of the answer]
\label{lem:IPhO2026Problems_problem_IPhO_2026_3_A_1:meanRadius_form_eq_volume_form}
\lean{IPhO2026.Problem3.PartA1.ParamagneticTorusA1.meanRadius_form_eq_volume_form}
\uses{lem:IPhO2026Problems_problem_IPhO_2026_3_A_1:mean_circumference_eq}
$N I / (2\pi R) = N I A / V$.
\end{lemma}
\begin{proof}
Substitute Bridge 3 into the denominator
(uses $2\pi R > 0$, $V > 0$).
\end{proof}

\begin{theorem}[A.1 main target: $H = NIA/V$]
\label{thm:IPhO2026Problems_problem_IPhO_2026_3_A_1:paramagneticTorus_H_eq}
\lean{IPhO2026.Problem3.PartA1.paramagneticTorus_H_eq}
\uses{lem:IPhO2026Problems_problem_IPhO_2026_3_A_1:fieldMagnitude_eq_meanRadius_form, lem:IPhO2026Problems_problem_IPhO_2026_3_A_1:meanRadius_form_eq_volume_form}
The magnitude of $\vec{H}$ inside the paramagnetic torus, in terms of
$N$, $A$, $V$ and the instantaneous current $I$, is
$H = N I A / V$.
\end{theorem}
\begin{proof}
Chain Bridge 2 ($H = NI/(2\pi R)$) with Bridge 4
($NI/(2\pi R) = NIA/V$).
\end{proof}

\begin{theorem}[A.1 answer in mean-radius form]
\label{thm:IPhO2026Problems_problem_IPhO_2026_3_A_1:paramagneticTorus_H_eq_meanRadius}
\lean{IPhO2026.Problem3.PartA1.paramagneticTorus_H_eq_meanRadius}
\uses{lem:IPhO2026Problems_problem_IPhO_2026_3_A_1:fieldMagnitude_eq_meanRadius_form}
Equivalent form obtained directly from Ampère's law before the geometry
substitution: $H = N I / (2\pi R)$.
\end{theorem}
\begin{proof}
Restatement of Bridge 2.
\end{proof}

% --- Archon physics formalization source end ---
